\ExerciseChapter{因子分析}

\label{chapter:2}

\section{往年题}

\par\addvspace{6pt}\noindent\begin{minipage}{\linewidth}

\subsection{单项选择题}

\Question{M15-3}{正交旋转的不变量}

对因子初始解作正交旋转，没有改变的是：


\begin{choices}

\item 原变量共同度。

\item 各个因子负荷。

\item 因子得分系数。

\item 单个共同因子的解释贡献。

\item 共同因子与原变量的各个相关系数。

\end{choices}

\end{minipage}\par\vspace{4pt}

\par\addvspace{6pt}\noindent\begin{minipage}{\linewidth}

\Question{M17-4}{探索性因子分析}

关于探索性因子分析，正确的是：


\begin{choices}

\item 给定数据的共同因子解是唯一的。

\item 因子旋转的目的就是获得因子间相关。

\item 因子命名主要依据负荷大小及变量含义。

\item 第一因子得分方差必不小于其余得分方差。

\item 变量相关较强通常不利于因子分析。

\end{choices}

\end{minipage}\par\vspace{4pt}

\par\addvspace{6pt}\noindent\begin{minipage}{\linewidth}

\Question{M21-4}{探索性因子分析}

关于探索性因子分析，正确的是：


\begin{choices}

\item 给定数据的共同因子解是唯一的。

\item 因子旋转的目的就是获得因子间相关。

\item 因子命名主要依据负荷大小及变量含义。

\item 第一因子得分方差必不小于其余得分方差。

\item 变量相关较强通常不利于因子分析。

\end{choices}

\end{minipage}\par\vspace{4pt}

\par\addvspace{6pt}\noindent\begin{minipage}{\linewidth}

\Question{G17-3}{同一学生在不同指标下的排序}

\begin{center}\small\setlength{\tabcolsep}{4pt}\renewcommand{\arraystretch}{1.18}\begin{tabular}{@{}lrrr@{}}
\toprule
姓名 & 语文 & 数学 & 英语 \\
\midrule

马云 & 84 & 88 & 80 \\
方舟子 & 95 & 99 & 95 \\
三毛 & 98 & 20 & 90 \\
王宝强 & 60 & 45 & 55 \\
赵薇 & 85 & 50 & 80 \\
\bottomrule\end{tabular}\end{center}

数据作Z变换，以主成分法提取2因子并作varimax旋转。对马云较有利的综合评价指标是：I第一因子得分；II第二因子得分；III第一主成分得分；IV三科总分。


\begin{choices}

\item 仅I。

\item 仅II。

\item I与III。

\item II与IV。

\item II与III。

\end{choices}

\end{minipage}\par\vspace{4pt}

\par\addvspace{6pt}\noindent\begin{minipage}{\linewidth}

\Question{M19-1}{因子提取与方差溢出}

对正半定相关矩阵提取少于或等于全部成分时，下列哪种提取方式保证其提取共同度不超过1，从而不会出现共同度大于1的Heywood现象？


\begin{choices}

\item 主成分提取法。

\item 不加约束的迭代主轴因子法。

\item 不加约束的极大似然因子法。

\item 任意指定负独特方差的模型。

\item 把共同度初值全部设为2的方法。

\end{choices}

\end{minipage}\par\vspace{4pt}

\par\addvspace{6pt}\noindent\begin{minipage}{\linewidth}

\Question{M19-2}{因子模型中的系数}

写标准化变量的因子模型 \(Z_j=\sum_k l_{jk}F_k+\varepsilon_j\) 时，\(l_{jk}\) 应从哪项输出读取？


\begin{choices}

\item 因子负荷矩阵（指定旋转阶段）。

\item 因子得分系数矩阵。

\item 各样本的因子得分矩阵。

\item 原始观测数据矩阵。

\item KMO统计量。

\end{choices}

\end{minipage}\par\vspace{4pt}

\par\addvspace{6pt}\noindent\begin{minipage}{\linewidth}

\Question{M19-3}{保留因子数的基本考虑}

选择因子数时，最合理的总体原则是：


\begin{choices}

\item 无论资料如何都保留2因子。

\item 只要第一特征值大于1就不再看其他信息。

\item 用尽可能简洁且可解释的模型保留足够信息，并结合拟合和专业目的判断。

\item 只比较变量原始均数。

\item 因子数必须等于样本量。

\end{choices}

\end{minipage}\par\vspace{4pt}

\par\addvspace{6pt}\noindent\begin{minipage}{\linewidth}

\Question{N19-7}{正交因子共同度}

两个标准正交因子对某变量的负荷为0.6和0.3，共同度为：


\begin{choices}

\item 0.90。

\item 0.45。

\item 0.18。

\item 1.80。

\item 0.30。

\end{choices}

\end{minipage}\par\vspace{4pt}

\par\addvspace{6pt}\noindent\begin{minipage}{\linewidth}

\Question{N19-8}{斜交旋转的模式与结构}

斜交因子模型模式矩阵为 \(P\)、因子相关矩阵为 \(\Phi\)，结构矩阵为：


\begin{choices}

\item \(P\)必定不变。

\item \(P\Phi\)。

\item \(P+\Phi\)。

\item \(P^TP\)。

\item \(\Phi^{-1}\)。

\end{choices}

\end{minipage}\par\vspace{4pt}

\par\addvspace{6pt}\noindent\begin{minipage}{\linewidth}

\Question{N21-12}{非迭代SMC初值}

相关阵可逆。主轴因子法的SMC初始共同度应取：


\begin{choices}

\item \(1-1/(R^{-1})_{jj}\)。

\item \((R^{-1})_{jj}\)。

\item 任意大于1的值。

\item 原变量均值。

\item 全部取0且不再提取。

\end{choices}

\end{minipage}\par\vspace{4pt}

\par\addvspace{6pt}\noindent\begin{minipage}{\linewidth}

\Question{N21-13}{因子评分方法的识别}

下式属于何种因子评分？\(\widehat f=(L^T\Psi^{-1}L)^{-1}L^T\Psi^{-1}z\)。


\begin{choices}

\item Bartlett加权最小二乘评分。

\item Thomson回归评分。

\item KMO计算。

\item Ward更新公式。

\item Bayes后验概率。

\end{choices}

\end{minipage}\par\vspace{4pt}

\par\addvspace{6pt}\noindent\begin{minipage}{\linewidth}

\Question{N21-14}{KMO的含义}

KMO较大通常意味着：


\begin{choices}

\item 偏相关相对于简单相关较小，资料更适于探索共享因子结构。

\item 所有变量均独立。

\item 模型已被证明正确。

\item 因子数已经确定。

\item 所有因子得分完全无误差。

\end{choices}

\end{minipage}\par\vspace{4pt}

\par\addvspace{6pt}\noindent\begin{minipage}{\linewidth}

\subsection{简答题}

\Question{FS20}{因子负荷的行、列平方和}

什么是因子负荷？其行平方和与列平方和各代表什么含义？


\worklines{3}

\end{minipage}\par\vspace{4pt}

\subsection{计算与分析题}

\Question{F13}{戒毒者量表：初始解与旋转解}

调查20名戒毒者的抑郁 \(X_1\)、焦虑 \(X_2\)、情绪管理 \(X_3\)、情绪判读 \(X_4\)、与他人关系 \(X_5\)、社会满意度 \(X_6\)、工作满意度 \(X_7\)。用主成分法提取因子，对初始模型作方差最大正交旋转。共20分。 1. 写出``焦虑（\(X_3\)）''的共同度并说明含义（原题名称与编号矛盾）。 2. 写出第四因子的贡献率并说明含义。 3. 写出累计贡献率并说明含义。 4. 写出焦虑 \(X_2\) 的因子模型及它与第一因子的相关系数。 5. 写出第三因子得分式，求其均值、标准差及与第一因子的简单相关系数。

\begin{center}\small\setlength{\tabcolsep}{4pt}\renewcommand{\arraystretch}{1.18}\begin{tabular}{@{}lrr@{}}
\toprule
变量 & 初始共同度 & 提取共同度 \\
\midrule

\(X_1\) & 1 & 0.911 \\
\(X_2\) & 1 & 0.902 \\
\(X_3\) & 1 & 0.803 \\
\(X_4\) & 1 & 0.723 \\
\(X_5\) & 1 & 0.72 \\
\(X_6\) & 1 & 0.913 \\
\(X_7\) & 1 & 0.896 \\
\bottomrule\end{tabular}\end{center}

初始与旋转后的方差解释表。

\begin{center}\small\setlength{\tabcolsep}{4pt}\renewcommand{\arraystretch}{1.18}\begin{tabular}{@{}lrrrrrr@{}}
\toprule
因子 & 初始\(\lambda\) & 贡献\% & 累计\% & 旋转\(\lambda\) & 贡献\% & 累计\% \\
\midrule

1 & 2.053 & 29.327 & 29.327 & 1.997 & 28.523 & 28.523 \\
2 & 1.695 & 24.22 & 53.547 & 1.499 & 21.413 & 49.936 \\
3 & 1.1 & 15.708 & 69.255 & 1.234 & 17.633 & 67.569 \\
4 & 1.02 & 14.577 & 83.832 & 1.138 & 16.263 & 83.832 \\
5 & 0.639 & 9.124 & 92.956 & & & \\
6 & 0.385 & 5.494 & 98.45 & & & \\
7 & 0.108 & 1.55 & 100 & & & \\
\bottomrule\end{tabular}\end{center}

初始载荷矩阵。

\begin{center}\small\setlength{\tabcolsep}{4pt}\renewcommand{\arraystretch}{1.18}\begin{tabular}{@{}lrrrr@{}}
\toprule
变量 & 1 & 2 & 3 & 4 \\
\midrule

\(X_1\) & 0.902 & -0.177 & 0.139 & -0.217 \\
\(X_2\) & 0.886 & -0.213 & 0.135 & 0.228 \\
\(X_3\) & 0.078 & 0.718 & 0.412 & -0.334 \\
\(X_4\) & 0.203 & 0.774 & 0.088 & -0.273 \\
\(X_5\) & 0.587 & 0.351 & -0.258 & 0.431 \\
\(X_6\) & -0.239 & 0.426 & 0.355 & 0.741 \\
\(X_7\) & -0.07 & -0.446 & 0.832 & 0.011 \\
\bottomrule\end{tabular}\end{center}

旋转载荷矩阵。

\begin{center}\small\setlength{\tabcolsep}{4pt}\renewcommand{\arraystretch}{1.18}\begin{tabular}{@{}lrrrr@{}}
\toprule
变量 & 1 & 2 & 3 & 4 \\
\midrule

\(X_1\) & 0.878 & 0.13 & 0.082 & -0.341 \\
\(X_2\) & 0.943 & -0.098 & 0.028 & 0.034 \\
\(X_3\) & -0.024 & 0.887 & 0.089 & 0.085 \\
\(X_4\) & 0.052 & 0.81 & -0.248 & 0.041 \\
\(X_5\) & 0.553 & 0.09 & -0.54 & 0.338 \\
\(X_6\) & -0.114 & 0.111 & 0.053 & 0.941 \\
\(X_7\) & 0.116 & -0.086 & 0.929 & 0.111 \\
\bottomrule\end{tabular}\end{center}

旋转后的成分得分系数矩阵（作用于标准化变量）。

\begin{center}\small\setlength{\tabcolsep}{4pt}\renewcommand{\arraystretch}{1.18}\begin{tabular}{@{}lrrrr@{}}
\toprule
变量 & 1 & 2 & 3 & 4 \\
\midrule

\(X_1\) & 0.421 & 0.116 & 0.115 & -0.248 \\
\(X_2\) & 0.49 & -0.095 & 0.056 & 0.124 \\
\(X_3\) & -0.025 & 0.626 & 0.189 & -0.022 \\
\(X_4\) & -0.003 & 0.532 & -0.102 & -0.057 \\
\(X_5\) & 0.285 & -0.059 & -0.407 & 0.329 \\
\(X_6\) & 0.021 & -0.018 & 0.082 & 0.837 \\
\(X_7\) & 0.114 & 0.042 & 0.778 & 0.149 \\
\bottomrule\end{tabular}\end{center}


\worklines{3}

\Question{FU}{戒毒者量表：旋转后模型}

调查20名戒毒者的抑郁 \(X_1\)、焦虑 \(X_2\)、情绪管理 \(X_3\)、情绪判读 \(X_4\)、与他人关系 \(X_5\)、社会满意度 \(X_6\)、工作满意度 \(X_7\)。用主成分法提取因子，对初始模型作方差最大正交旋转。共20分。 1. 写出焦虑 \(X_2\) 的共同度及其含义。 2. 写出旋转后第三因子的贡献率及其含义。 3. 写出总累计贡献率及其含义。 4. 写出焦虑 \(X_2\) 的因子模型及其与第一因子的相关系数。 5. 写出旋转后第三因子得分式，求其均值、标准差及与第一因子得分的简单相关系数。

\begin{center}\small\setlength{\tabcolsep}{4pt}\renewcommand{\arraystretch}{1.18}\begin{tabular}{@{}lrr@{}}
\toprule
变量 & 初始共同度 & 提取共同度 \\
\midrule

\(X_1\) & 1 & 0.911 \\
\(X_2\) & 1 & 0.902 \\
\(X_3\) & 1 & 0.803 \\
\(X_4\) & 1 & 0.723 \\
\(X_5\) & 1 & 0.72 \\
\(X_6\) & 1 & 0.913 \\
\(X_7\) & 1 & 0.896 \\
\bottomrule\end{tabular}\end{center}

初始与旋转后的方差解释表。

\begin{center}\small\setlength{\tabcolsep}{4pt}\renewcommand{\arraystretch}{1.18}\begin{tabular}{@{}lrrrrrr@{}}
\toprule
因子 & 初始\(\lambda\) & 贡献\% & 累计\% & 旋转\(\lambda\) & 贡献\% & 累计\% \\
\midrule

1 & 2.053 & 29.327 & 29.327 & 1.997 & 28.523 & 28.523 \\
2 & 1.695 & 24.22 & 53.547 & 1.499 & 21.413 & 49.936 \\
3 & 1.1 & 15.708 & 69.255 & 1.234 & 17.633 & 67.569 \\
4 & 1.02 & 14.577 & 83.832 & 1.138 & 16.263 & 83.832 \\
5 & 0.639 & 9.124 & 92.956 & & & \\
6 & 0.385 & 5.494 & 98.45 & & & \\
7 & 0.108 & 1.55 & 100 & & & \\
\bottomrule\end{tabular}\end{center}

初始载荷矩阵。

\begin{center}\small\setlength{\tabcolsep}{4pt}\renewcommand{\arraystretch}{1.18}\begin{tabular}{@{}lrrrr@{}}
\toprule
变量 & 1 & 2 & 3 & 4 \\
\midrule

\(X_1\) & 0.902 & -0.177 & 0.139 & -0.217 \\
\(X_2\) & 0.886 & -0.213 & 0.135 & 0.228 \\
\(X_3\) & 0.078 & 0.718 & 0.412 & -0.334 \\
\(X_4\) & 0.203 & 0.774 & 0.088 & -0.273 \\
\(X_5\) & 0.587 & 0.351 & -0.258 & 0.431 \\
\(X_6\) & -0.239 & 0.426 & 0.355 & 0.741 \\
\(X_7\) & -0.07 & -0.446 & 0.832 & 0.011 \\
\bottomrule\end{tabular}\end{center}

旋转载荷矩阵。

\begin{center}\small\setlength{\tabcolsep}{4pt}\renewcommand{\arraystretch}{1.18}\begin{tabular}{@{}lrrrr@{}}
\toprule
变量 & 1 & 2 & 3 & 4 \\
\midrule

\(X_1\) & 0.878 & 0.13 & 0.082 & -0.341 \\
\(X_2\) & 0.943 & -0.098 & 0.028 & 0.034 \\
\(X_3\) & -0.024 & 0.887 & 0.089 & 0.085 \\
\(X_4\) & 0.052 & 0.81 & -0.248 & 0.041 \\
\(X_5\) & 0.553 & 0.09 & -0.54 & 0.338 \\
\(X_6\) & -0.114 & 0.111 & 0.053 & 0.941 \\
\(X_7\) & 0.116 & -0.086 & 0.929 & 0.111 \\
\bottomrule\end{tabular}\end{center}

旋转后的成分得分系数矩阵（作用于标准化变量）。

\begin{center}\small\setlength{\tabcolsep}{4pt}\renewcommand{\arraystretch}{1.18}\begin{tabular}{@{}lrrrr@{}}
\toprule
变量 & 1 & 2 & 3 & 4 \\
\midrule

\(X_1\) & 0.421 & 0.116 & 0.115 & -0.248 \\
\(X_2\) & 0.49 & -0.095 & 0.056 & 0.124 \\
\(X_3\) & -0.025 & 0.626 & 0.189 & -0.022 \\
\(X_4\) & -0.003 & 0.532 & -0.102 & -0.057 \\
\(X_5\) & 0.285 & -0.059 & -0.407 & 0.329 \\
\(X_6\) & 0.021 & -0.018 & 0.082 & 0.837 \\
\(X_7\) & 0.114 & 0.042 & 0.778 & 0.149 \\
\bottomrule\end{tabular}\end{center}


\worklines{3}

\Question{F15}{三科成绩因子模型与得分}

下列为三科考试成绩的因子分析结果。所有初始共同度为1；以主成分法提取，采用Kaiser标准化的varimax旋转，3次迭代收敛。共15分。

\begin{enumerate}
\def\labelenumi{\arabic{enumi}.}
\tightlist
\item
  分析者选择了几个公因子？
\item
  语文、数学的共同度各是多少？
\item
  各公因子的方差贡献率各是多少？
\item
  列出因子模型方程。
\item
  列出因子得分模型方程。
\end{enumerate}

\begin{center}\small\setlength{\tabcolsep}{4pt}\renewcommand{\arraystretch}{1.18}\begin{tabular}{@{}lrrr@{}}
\toprule
变量 & 均值 & 标准差 & 共同度 \\
\midrule

语文 & 82.2 & 11.946 & 0.978 \\
数学 & 75.2 & 19.942 & 1 \\
英语 & 80.2 & 7.95 & 0.979 \\
\bottomrule\end{tabular}\end{center}

\begin{center}\small\setlength{\tabcolsep}{4pt}\renewcommand{\arraystretch}{1.18}\begin{tabular}{@{}lrrrrrr@{}}
\toprule
成分 & 初始\(\lambda\) & 贡献\% & 累计\% & 旋转\(\lambda\) & 贡献\% & 累计\% \\
\midrule

1 & 2.223 & 74.109 & 74.109 & 1.912 & 63.736 & 63.736 \\
2 & 0.734 & 24.462 & 98.57 & 1.045 & 34.835 & 98.57 \\
3 & 0.043 & 1.43 & 100 & & & \\
\bottomrule\end{tabular}\end{center}

\begin{center}\small\setlength{\tabcolsep}{4pt}\renewcommand{\arraystretch}{1.18}\begin{tabular}{@{}lrrrrrr@{}}
\toprule
变量 & 初始1 & 初始2 & 旋转1 & 旋转2 & 得分1 & 得分2 \\
\midrule

语文 & 0.963 & -0.225 & 0.959 & 0.24 & 0.525 & -0.075 \\
数学 & 0.632 & 0.775 & 0.208 & 0.978 & -0.229 & 1.069 \\
英语 & 0.947 & -0.289 & 0.974 & 0.176 & 0.558 & -0.155 \\
\bottomrule\end{tabular}\end{center}


\worklines{3}

\par\addvspace{6pt}\noindent\begin{minipage}{\linewidth}

\Question{F17}{成绩分析：旋转载荷与因子命名}

对下表进行主成分与因子分析，所有变量作Z变换，以主成分法提取2个公因子并作varimax旋转。

\begin{center}\small\setlength{\tabcolsep}{4pt}\renewcommand{\arraystretch}{1.18}\begin{tabular}{@{}lrrr@{}}
\toprule
姓名 & 语文 & 数学 & 英语 \\
\midrule

马云 & 84 & 88 & 80 \\
方舟子 & 95 & 99 & 95 \\
三毛 & 98 & 20 & 90 \\
王宝强 & 60 & 45 & 55 \\
赵薇 & 85 & 50 & 80 \\
\bottomrule\end{tabular}\end{center}

\begin{enumerate}
\def\labelenumi{\arabic{enumi}.}
\tightlist
\item
  写出因子模型方程。（Q1，5分）
\item
  给第1、第2因子命名。（Q2，5分）
\end{enumerate}


\worklines{3}

\end{minipage}\par\vspace{4pt}

\par\addvspace{6pt}\noindent\begin{minipage}{\linewidth}

\Question{F19}{医院满意度：主成分提取与Bartlett得分}

对基层医院医护人员满意度数据（223人、17项指标）进行因子分析：用主成分法提取4个因子，作正交旋转，用Bartlett法计算因子得分。共20分。变量如下。

\begin{center}\small\setlength{\tabcolsep}{4pt}\renewcommand{\arraystretch}{1.18}\begin{tabular}{@{}lrrr@{}}
\toprule
变量 & 含义 & 变量 & 含义 \\
\midrule

\(X_1\) & 业务用房 & \(X_10\) & 待遇公平 \\
\(X_2\) & 医疗设备 & \(X_11\) & 职工利益 \\
\(X_3\) & 技术水平 & \(X_12\) & 管理水平 \\
\(X_4\) & 办公条件 & \(X_13\) & 听取意见 \\
\(X_5\) & 报酬收入 & \(X_14\) & 社会尊重 \\
\(X_6\) & 发展信心 & \(X_15\) & 发展前途 \\
\(X_7\) & 住房分配 & \(X_16\) & 适应工作 \\
\(X_8\) & 住房现状 & \(X_17\) & 福利待遇 \\
\(X_9\) & 考评制度 & & \\
\bottomrule\end{tabular}\end{center}

\begin{enumerate}
\def\labelenumi{\arabic{enumi}.}
\tightlist
\item
  写出 \(X_1\) 的因子表达式。
\item
  参考各变量含义，给4个因子命名。
\item
  给出case100的第二因子得分；若数据来自抽样调查，估计全人群第二因子维度满意度水平。
\end{enumerate}


\worklines{3}

\end{minipage}\par\vspace{4pt}

\par\addvspace{6pt}\noindent\begin{minipage}{\linewidth}

\Question{F20}{医院满意度：非迭代SMC与正交旋转}

使用医院满意度17项数据做因子分析。共20分。

\begin{enumerate}
\def\labelenumi{\arabic{enumi}.}
\tightlist
\item
  计算KMO及各变量MSA（原稿写KMO/SMA），判断适合度。（4分）
\item
  用主因子法中的SMC法，不迭代，提取4个共同因子，作varimax旋转。列出累计特征值、共同度和负荷矩阵。（8分）
\item
  依据第2问结果命名因子。（4分）
\item
  计算该4因子模型的Thomson回归法得分，列出前5个case的得分。（4分）
\end{enumerate}

变量含义与医院满意度数据表一致；数据文件已附。


\worklines{3}

\end{minipage}\par\vspace{4pt}

\par\addvspace{6pt}\noindent\begin{minipage}{\linewidth}

\Question{F21}{医院满意度：Promax斜交旋转}

用主因子法的SMC法（不迭代）提取4个因子，作promax斜交旋转，幂指数为4。

\begin{enumerate}
\def\labelenumi{\arabic{enumi}.}
\tightlist
\item
  按负荷与变量含义给4个共同因子命名。（4分）
\item
  填空（每空1分）：

  \begin{itemize}
  \tightlist
  \item
    4因子的累计特征值（以初始解或正交旋转解为准）。
  \item
    \(X_1\) 的共同度。
  \item
    \(X_2\) 在第二因子上的负荷。
  \item
    第一与第四因子的相关系数。
  \item
    case2在第一因子上的Bartlett得分。
  \item
    case1在第一因子上的Thomson回归得分。
  \end{itemize}
\end{enumerate}


\worklines{3}

\end{minipage}\par\vspace{4pt}

\par\addvspace{6pt}\noindent\begin{minipage}{\linewidth}

\subsection{文献分析题}

\Question{L13-3}{文献中的两因子模型}

所指原始文献、图1及完整数据没有随题卷提供。以下保留设问；不能据此独立核验文献中的具体研究事实。 1. 8项指标得到两因子模型时，哪项共同度最小？ 2. 分析作者首先取两因子模型的可能理由。


\worklines{3}

\end{minipage}\par\vspace{4pt}

\par\addvspace{6pt}\noindent\begin{minipage}{\linewidth}

\Question{L15}{文献复现与方法评价}

请仔细阅读给出的文献，重复文献的分析结果，并作适当评价或改进。原题35分。

文献及其数据未附。


\worklines{3}

\end{minipage}\par\vspace{4pt}

\par\addvspace{6pt}\noindent\begin{minipage}{\linewidth}

\Question{LU}{文献述评}

请阅读所给文献，并作述评。原题25分。

文献未附。


\worklines{3}

\end{minipage}\par\vspace{4pt}

\clearpage\section{补充练习}

本节为配合正文新编的补充练习。除注明沿用正文例题外，数值与研究摘要均为教学设定。题型与作答方式沿用本章往年题，解析见章末。

\par\addvspace{6pt}\noindent\begin{minipage}{\linewidth}

\subsection{单项选择题}

\Question{SUP-F1}{最大似然因子模型的自由度}

对6个指标拟合正交公因子模型。按正文公式 \(df=[(p-m)^2-p-m]/2\)，其中 \(p\) 为指标数，\(m\) 为公因子数。以下正确的是：


\begin{choices}

\item 两因子模型的自由度为6。

\item 三因子模型的自由度为1。

\item 两因子模型自由度为4；三因子模型自由度为0，后者不能按通常的正自由度卡方检验评价拟合。

\item 公因子数越多，模型自由度越大。

\item 自由度为0表明模型已被证明正确。

\end{choices}

\end{minipage}\par\vspace{4pt}

\par\addvspace{6pt}\noindent\begin{minipage}{\linewidth}

\subsection{简答题}

\Question{SUP-F2}{三种因子提取方法的区别}

从相关矩阵出发时，主成分法、主因子法和最大似然法分别怎样获得初始因子解？说明主成分法与主因子法对相关矩阵对角线的处理差别，以及正文所述最大似然因子分析进行拟合优度检验的主要条件。


\worklines{3}

\end{minipage}\par\vspace{4pt}

\par\addvspace{6pt}\noindent\begin{minipage}{\linewidth}

\subsection{计算与分析题}

\Question{SUP-F3}{体格指标的共同度与还原相关}

以收缩压、舒张压、体重、腰围4个标准化指标拟合两因子模型。公因子均值为0、方差为1且互不相关，特殊因子与公因子不相关，特殊因子之间也不相关。某次分析给出的正交旋转载荷如下（教学设定）。

\begin{center}\small\setlength{\tabcolsep}{4pt}\renewcommand{\arraystretch}{1.18}\begin{tabular}{@{}lrr@{}}
\toprule
变量 & 因子1 & 因子2 \\
\midrule

收缩压 & 0.8 & 0.2 \\
舒张压 & 0.7 & 0.1 \\
体重 & 0.1 & 0.8 \\
腰围 & 0.3 & 0.7 \\
\bottomrule\end{tabular}\end{center}

共15分。

\begin{enumerate}
\def\labelenumi{\arabic{enumi}.}
\tightlist
\item
  写出收缩压的因子模型，并计算4个变量的共同度和特殊方差。（5分）
\item
  计算两个因子各自的贡献率及累计贡献率，给因子命名。（4分）
\item
  计算模型还原的收缩压与舒张压相关系数。若样本相关系数为0.62，残差相关（样本值减模型还原值）是多少？（4分）
\item
  能否将载荷表的第一列直接乘以某人的4项标准分，称为其Thomson回归法因子得分？（2分）
\end{enumerate}


\worklines{3}

\end{minipage}\par\vspace{4pt}

\par\addvspace{6pt}\noindent\begin{minipage}{\linewidth}

\subsection{文献分析题}

\Question{SUP-F4}{两因子模型的研究结果评议}

以下为依据课堂方法编写的模拟研究摘要，全部评议所需结果已列出，不对应真实论文。

研究者测量220名受试者的6项数值指标，认为独立观测及多元正态等分析条件可以接受。最大似然提取1个因子时，拟合检验 \(\chi^2=28.4\)，\(df=9\)，\(P<0.001\)；提取2个因子时，\(\chi^2=3.6\)，\(df=4\)，\(P=0.463\)。两因子模型旋转前后列平方和如下。

\begin{center}\small\setlength{\tabcolsep}{4pt}\renewcommand{\arraystretch}{1.18}\begin{tabular}{@{}lrr@{}}
\toprule
阶段 & 因子1 & 因子2 \\
\midrule

旋转前 & 2.7 & 1.2 \\
varimax旋转后 & 2.05 & 1.85 \\
\bottomrule\end{tabular}\end{center}

作者写道：``两因子检验不显著，证明真实结构恰好只有两个因子。旋转后第二因子贡献增加，说明旋转提高了模型的总解释率。''共15分。

\begin{enumerate}
\def\labelenumi{\arabic{enumi}.}
\tightlist
\item
  写出拟合优度检验的零假设，在 \(\alpha=0.05\) 下比较这两个模型。（5分）
\item
  计算旋转前后各因子及累计贡献率，评价作者关于旋转的结论。（5分）
\item
  修正``证明真实结构恰好只有两个因子''的表述，并说明因子命名还需什么信息。（5分）
\end{enumerate}


\worklines{3}

\end{minipage}\par\vspace{4pt}
\clearpage\section{习题解析}

\begin{keyidea}[核对方式]按题号核对。补写题目、原稿歧义及复算约定列在对应解析中。\end{keyidea}

\Answer{M15-3}{正交旋转的不变量}

\textbf{答案：A。}

若 \(L^*=LT\)、\(TT^T=I\)，则 \(L^*L^{*T}=LL^T\)，共同度不变；负荷和单列解释贡献一般会重新分配。


\Answer{M17-4}{探索性因子分析}

\textbf{答案：C。}

旋转追求更易解释的简单结构；正交旋转保持因子不相关，斜交旋转允许相关，但相关不是唯一目的。因子命名依赖负荷模式与测量内容，不能只机械选最大数值。


\Answer{M21-4}{探索性因子分析}

\textbf{答案：C。}

旋转追求更易解释的简单结构；正交旋转保持因子不相关，斜交旋转允许相关，但相关不是唯一目的。因子命名依赖负荷模式与测量内容，不能只机械选最大数值。


\Answer{G17-3}{同一学生在不同指标下的排序}

\textbf{答案：D。}

按语文英语为第一因子、数学为第二因子，马云的第二因子和总分均位列第2，第一因子较靠后，第一主成分约第3。判断``有利''必须与因子含义和排序相联系，不能只比较不同尺度的得分绝对值。


\Answer{M19-1}{因子提取与方差溢出}

\textbf{答案：A。}

PCA保留的载荷行平方和不超过全部成分行平方和，而后者为相关阵对角元1。该结论针对有效正半定相关阵；缺失处理产生非正定矩阵等情况仍可能导致其他计算问题。


\textbf{补写说明。} 2019回忆稿只记录主题，没有完整题干与选项；本题依据该主题补写，不宣称恢复原选项。


\Answer{M19-2}{因子模型中的系数}

\textbf{答案：A。}

变量由因子表示使用负荷；估计个体因子得分使用得分系数，两者方向不同。斜交时应明确模式负荷与结构相关。


\textbf{补写说明。} 2019回忆稿只记录主题，没有完整题干与选项；本题依据该主题补写，不宣称恢复原选项。


\Answer{M19-3}{保留因子数的基本考虑}

\textbf{答案：C。}

累计方差贡献率、碎石图、平行分析、模型拟合和可解释性均可作为依据。70\%或80\%属于经验阈值，不是所有数据都适用的根本定律。


\textbf{补写说明。} 2019回忆稿只记录主题，没有完整题干与选项；本题依据该主题补写，不宣称恢复原选项。


\Answer{N19-7}{正交因子共同度}

\textbf{答案：B。}

共同度为 \(0.6^2+0.3^2=0.45\)，独特方差为0.55。


\textbf{补写说明。} 2019卷该部分仅有总题数和部分回忆，整题内容未记录。本题为按课程知识点新编的补充练习，不是恢复的往年原题；编号只用于本包覆盖核对。


\Answer{N19-8}{斜交旋转的模式与结构}

\textbf{答案：B。}

结构矩阵给变量与因子的相关；因子相关不为单位阵时，通常不同于模式矩阵。


\textbf{补写说明。} 2019卷该部分仅有总题数和部分回忆，整题内容未记录。本题为按课程知识点新编的补充练习，不是恢复的往年原题；编号只用于本包覆盖核对。


\Answer{N21-12}{非迭代SMC初值}

\textbf{答案：A。}

SMC是第j个变量被其余变量线性解释的复相关系数平方。


\textbf{补写说明。} 2021卷该部分仅有总题数和部分回忆，整题内容未记录。本题为按课程知识点新编的补充练习，不是恢复的往年原题；编号只用于本包覆盖核对。


\Answer{N21-13}{因子评分方法的识别}

\textbf{答案：A。}

该式对独特误差方差加权，最小化给定共同因子模型中的加权残差平方和。


\textbf{补写说明。} 2021卷该部分仅有总题数和部分回忆，整题内容未记录。本题为按课程知识点新编的补充练习，不是恢复的往年原题；编号只用于本包覆盖核对。


\Answer{N21-14}{KMO的含义}

\textbf{答案：A。}

KMO比较简单相关平方与偏相关平方的相对大小；它是适合度诊断之一，不替代模型与测量论证。


\textbf{补写说明。} 2021卷该部分仅有总题数和部分回忆，整题内容未记录。本题为按课程知识点新编的补充练习，不是恢复的往年原题；编号只用于本包覆盖核对。


\Answer{FS20}{因子负荷的行、列平方和}

因子负荷是共同因子模型中原始变量对共同因子的系数。标准化变量、单位方差且正交的因子下，负荷等于变量与因子的相关系数；第 \(j\) 行平方和为变量共同度 \(h_j^2\)，第 \(k\) 列平方和为该因子对全部标准化变量的解释方差，除以变量数得到贡献率。斜交旋转时，模式矩阵 \(P\) 一般不等于相关系数矩阵；结构矩阵为 \(P\Phi\)，共同度应取 \(\operatorname{diag}(P\Phi P^T)\)。


\Answer{F13}{戒毒者量表：初始解与旋转解}

令 \(Z_j=(X_j-\bar X_j)/s_j\)。共同度是所保留因子对某个标准化变量方差的解释比例。\(X_2\)为0.902、\(X_3\)为0.803。

初始第四因子贡献率为14.577\%；旋转后第三因子为17.633\%。各题要求的旋转阶段不能混用。总累计贡献率为83.832\%，正交旋转不改变这个总量。

焦虑的旋转模型为 \[Z_2=.943F_1-.098F_2+.028F_3+.034F_4+\varepsilon_2,\] 其中独特部分方差约 \(1-.902=.098\)；共同因子标准化且正交时，\(\operatorname{Corr}(Z_2,F_1)=.943\)。若用初始解，则相应系数为(.886,−.213, .135, .228)，与第一初始因子相关为.886。

第三得分表达式： \[\begin{aligned}\widehat F_3={}&.115Z_1+.056Z_2+.189Z_3-.102Z_4\\&-.407Z_5+.082Z_6+.778Z_7.\end{aligned}\] 以本题主成分提取并标准化成分得分的正交旋转约定，其样本均值约0、标准差约1，与第一得分相关约0；舍入后的系数会引起小偏差。一般共同因子模型的估计得分未必恰有单位方差或零相关，不能不加条件推广。

\textbf{原题校正。} 第1问按变量编号 \(X_3\) 作答为0.803，但其正确名称为情绪管理；若按名称``焦虑''即 \(X_2\) 作答应为0.902。第2问未注明旋转后，原参考解按初始第四因子14.577\%作答；旋转后第四因子实际为16.263\%。


\Answer{FU}{戒毒者量表：旋转后模型}

令 \(Z_j=(X_j-\bar X_j)/s_j\)。共同度是所保留因子对某个标准化变量方差的解释比例。\(X_2\)为0.902、\(X_3\)为0.803。

初始第四因子贡献率为14.577\%；旋转后第三因子为17.633\%。各题要求的旋转阶段不能混用。总累计贡献率为83.832\%，正交旋转不改变这个总量。

焦虑的旋转模型为 \[Z_2=.943F_1-.098F_2+.028F_3+.034F_4+\varepsilon_2,\] 其中独特部分方差约 \(1-.902=.098\)；共同因子标准化且正交时，\(\operatorname{Corr}(Z_2,F_1)=.943\)。若用初始解，则相应系数为(.886,−.213, .135, .228)，与第一初始因子相关为.886。

第三得分表达式： \[\begin{aligned}\widehat F_3={}&.115Z_1+.056Z_2+.189Z_3-.102Z_4\\&-.407Z_5+.082Z_6+.778Z_7.\end{aligned}\] 以本题主成分提取并标准化成分得分的正交旋转约定，其样本均值约0、标准差约1，与第一得分相关约0；舍入后的系数会引起小偏差。一般共同因子模型的估计得分未必恰有单位方差或零相关，不能不加条件推广。


\Answer{F15}{三科成绩因子模型与得分}

保留2个因子。语文、数学共同度分别0.978、1.000。初始贡献为74.109\%、24.462\%；旋转后为63.736\%、34.835\%，累计均约98.570\%。

用标准化成绩 \(Z_1,Z_2,Z_3\) 表示旋转模型： \[\begin{aligned}Z_1&=.959F_1+.240F_2+\varepsilon_1,\\Z_2&=.208F_1+.978F_2+\varepsilon_2,\\Z_3&=.974F_1+.176F_2+\varepsilon_3.\end{aligned}\] 得分式为 \[\widehat F_1=.525Z_1-.229Z_2+.558Z_3,\] \[\widehat F_2=-.075Z_1+1.069Z_2-.155Z_3.\] 标准化必须使用题中各科均值和标准差。载荷矩阵用于``变量由因子表示''，得分系数矩阵用于``因子得分由变量估计''，不能交换，也不应再把旋转载荷任意除以其列平方和的平方根。


\Answer{F17}{成绩分析：旋转载荷与因子命名}

以``语文---英语''载荷为正统一因子方向，R复算旋转载荷如下：

\begin{center}\small\setlength{\tabcolsep}{4pt}\renewcommand{\arraystretch}{1.18}\begin{tabular}{@{}lrr@{}}
\toprule
变量 & 因子1 & 因子2 \\
\midrule

语文 & 0.9990 & 0.0232 \\
数学 & 0.0890 & 0.9960 \\
英语 & 0.9804 & 0.1931 \\
\bottomrule\end{tabular}\end{center}

令各科标准分为 \(Z\)，按每一行写 \(Z_j=l_{j1}F_1+l_{j2}F_2+\varepsilon_j\)。第一因子主要负荷在语文和英语，命名为``语文英语表现''；第二因子主要负荷在数学，命名为``数学表现''。因子顺序和符号可随软件约定整体改变，需连同载荷与得分一起调整。


\Answer{F19}{医院满意度：主成分提取与Bartlett得分}

采用varimax正交旋转，以负荷总体正向为符号约定： \[Z_1=.0760F_1+.8136F_2+.0045F_3+.2304F_4+\varepsilon_1.\] 四因子可依次概括为``管理与职工利益''\,``硬件与工作条件''\,``发展及适应''\,``收入住房福利''；具体命名应参照完整载荷矩阵（随包附出）。前三个/第四个初始特征值为7.0401、2.3927、1.3251、1.0165，累计解释约69.261\%。

固定载荷 \(L\) 和独特方差 \(\Psi\)，Bartlett估计为 \[\widehat f_B=(L^T\Psi^{-1}L)^{-1}L^T\Psi^{-1}z.\] 按附数据复算，case100第二因子得分约 \textbf{0.1331}。

若为简单随机样本且暂将评分规则视为固定，第二得分样本均值约0，标准差1.0032，SE约0.06718；均值的条件性95\%区间约 \([-0.1324, 0.1324]\)。该零点是样本标准化所设定的相对基准，不能据此说总体``满意度为0''。要估计有实际量尺的总体满意度，应明确抽样权重、评分标定和测量含义；复杂抽样须按设计估计标准误，并考虑评分规则估计的不确定性。


\Answer{F20}{医院满意度：非迭代SMC与正交旋转}

KMO约 \textbf{0.8988}；整体适合，但仍应结合设计、测量和相关结构。各变量MSA与载荷、共同度合并列在下表。

SMC以 \(h_j^{2(0)}=1-1/(R^{-1})_{jj}\) 替换相关阵对角线，进行一次特征分解，保留4个正特征根，不再更新共同度。4个特征值为6.630418、1.957903、0.810163、0.557710，和为 \textbf{9.956194}，相对17个标准化变量总方差的累计比例约 \textbf{58.566\%}。

varimax旋转载荷及共同度见下表。

\begin{center}\small\setlength{\tabcolsep}{4pt}\renewcommand{\arraystretch}{1.18}\begin{tabular}{@{}lrrrrrr@{}}
\toprule
变量 & 因子1 & 因子2 & 因子3 & 因子4 & 共同度 & MSA \\
\midrule

\(X_1\) & 0.0867 & 0.7507 & 0.0240 & 0.2064 & 0.6143 & 0.8638 \\
\(X_2\) & 0.1417 & 0.7552 & 0.0080 & 0.1776 & 0.6221 & 0.8705 \\
\(X_3\) & 0.2941 & 0.4654 & 0.3543 & 0.0825 & 0.4354 & 0.9099 \\
\(X_4\) & 0.1173 & 0.7789 & 0.0150 & 0.2052 & 0.6628 & 0.8513 \\
\(X_5\) & 0.3246 & 0.2126 & 0.0560 & 0.5385 & 0.4436 & 0.9116 \\
\(X_6\) & 0.4781 & 0.2034 & 0.4752 & 0.1282 & 0.5122 & 0.9226 \\
\(X_7\) & 0.2461 & 0.4211 & 0.0883 & 0.6176 & 0.6272 & 0.8998 \\
\(X_8\) & 0.1185 & 0.4205 & 0.1224 & 0.6168 & 0.5863 & 0.8841 \\
\(X_9\) & 0.7285 & 0.1748 & 0.1976 & 0.1386 & 0.6194 & 0.9046 \\
\(X_10\) & 0.7189 & 0.0795 & 0.2500 & 0.2656 & 0.6561 & 0.9414 \\
\(X_11\) & 0.7894 & 0.1962 & 0.1275 & 0.3012 & 0.7686 & 0.9313 \\
\(X_12\) & 0.7651 & 0.1748 & 0.2797 & 0.1178 & 0.7080 & 0.9157 \\
\(X_13\) & 0.7970 & 0.1175 & 0.1544 & 0.2382 & 0.7297 & 0.9149 \\
\(X_14\) & 0.3537 & -0.0878 & 0.5902 & 0.0541 & 0.4841 & 0.8496 \\
\(X_15\) & 0.4053 & 0.1370 & 0.5272 & 0.1418 & 0.4811 & 0.8895 \\
\(X_16\) & 0.0562 & 0.0259 & 0.6554 & 0.0649 & 0.4376 & 0.7173 \\
\(X_17\) & 0.4628 & 0.1031 & 0.1813 & 0.5568 & 0.5677 & 0.9156 \\
\bottomrule\end{tabular}\end{center}

因子可概括为管理/职工利益、硬件条件、发展与适应、待遇/住房福利。以样本相关阵作回归评分，得分系数为 \(R^{-1}L\)，得分矩阵为 \(Z R^{-1}L\)；前5人如下。

\begin{center}\small\setlength{\tabcolsep}{4pt}\renewcommand{\arraystretch}{1.18}\begin{tabular}{@{}lrrrr@{}}
\toprule
样本 & 因子1 & 因子2 & 因子3 & 因子4 \\
\midrule

1 & -0.7550 & -1.4696 & -0.1482 & -0.6744 \\
2 & -1.0415 & -0.7428 & -0.4462 & -0.8914 \\
3 & 0.0582 & 1.1280 & 0.0714 & 0.9176 \\
4 & -1.3451 & 0.6748 & 0.7096 & -0.5116 \\
5 & 0.0756 & 1.3890 & 0.8281 & -0.8834 \\
\bottomrule\end{tabular}\end{center}

\textbf{原解校正。} 原稿所列 \((L^T\Psi^{-1}L)^{-1}L^T\Psi^{-1}z\) 是Bartlett公式，不是Thomson回归公式。若使用拟合模型相关阵 \(LL^T+\Psi\) 替代样本 \(R\)，结果可能略有差异，必须标明。\texttt{factanal} 默认极大似然提取，不能代替题定的SMC不迭代提取。


\Answer{F21}{医院满意度：Promax斜交旋转}

以随包代码的因子次序与符号约定，命名同正交模型大致为管理/职工利益、硬件条件、发展/适应、待遇/住房福利。斜交因子之间有相关，负荷的解释应区分模式矩阵 \(P\) 与结构矩阵 \(P\Phi\)。

各空依次为：\textbf{9.956194；0.614270；0.776717；0.592298；−1.447396；−1.149144}。前者除以17约58.566\%。完整模式矩阵、结构矩阵、因子相关矩阵及全部得分随包提供。

共同度采用 \(h^2=\operatorname{diag}(P\Phi P^T)\)；Bartlett评分使用独特方差矩阵和模式矩阵；回归评分使用 \(Z R^{-1}P\Phi\)。

\textbf{原解差异。} 原参考稿的3.891是模式负荷矩阵某一列的平方和，不是四因子累计特征值；0.5936来自错误地对斜交模式矩阵直接作行平方和。其−1.4846、−1.1226使用了另一次极大似然提取后的得分，未忠实于SMC不迭代模型，故与这里的复算值不同。


\Answer{L13-3}{文献中的两因子模型}

原汇编在第1空写 \textbf{mn-sod}，但载荷表和文献未附，无法复核。正交因子下应计算各变量两载荷平方和，取最小者；斜交因子使用 \(P\Phi P^T\) 的对角元。

两因子选择可能根据特征根、碎石图、累计解释率、模型拟合和专业可解释性。应对照原文证据，不能把某个固定阈值当作者实际使用过的标准。


\Answer{L15}{文献复现与方法评价}

作答应包含研究问题与设计、测量和变量处理、所用多元方法、模型提取/选择步骤、结果复现，以及局限和改进。因子分析需核对提取与旋转方法、模式/结构矩阵、得分算法；预测模型需防止信息泄漏并进行验证。原文和数据缺失，本题不能完成数值复现；这不是补写题目选项能够解决的缺口。


\Answer{LU}{文献述评}

先明确研究目的、设计、人群和测量，再评估方法是否回答研究问题、假设与样本量是否合适、是否报告不确定性、变量筛选和验证是否规范。对结果解释区分关联、降维结构与预测能力。缺文献时只能给述评框架，不能编写文献实际结果。

\subsection{补充练习解析}

\Answer{SUP-F1}{最大似然因子模型的自由度}

\textbf{答案：C。}

两因子模型 \(df=[(6-2)^2-6-2]/2=4\)；三因子模型 \(df=[(6-3)^2-6-3]/2=0\)。增加公因子会减少剩余约束。正文的拟合优度检验要求模型自由度大于0；零自由度不能作为``模型正确''的证据。




\Answer{SUP-F2}{三种因子提取方法的区别}

\textbf{主成分法：} 对原相关矩阵（对角线为1）作特征分解，用所保留特征向量乘相应特征值的平方根组成载荷矩阵。

\textbf{主因子法：} 先估计各变量共同度，以共同度替换相关矩阵的对角线，得到约相关矩阵，再作特征分解。迭代主因子法会用新载荷重新估计共同度并继续更新；非迭代法只完成指定的一次提取。

\textbf{最大似然法：} 按共同因子模型 \(\Sigma=LL^T+\Psi\) 建立似然，估计载荷和特殊方差。正文采用多元正态模型；常规卡方拟合优度检验还要求独立观测、模型自由度大于0且估计满足通常的正则条件。若出现Heywood现象等异常估计，应先检查模型，不能机械解释检验结果。

因此，相关矩阵对角线均为1不能作为``主因子法已扣除特殊方差''的证据；软件默认的提取方法也必须与题意核对。




\Answer{SUP-F3}{体格指标的共同度与还原相关}

\textbf{（1）模型与方差分解。} 令 \(Z_1\) 为标准化收缩压，则 \[Z_1=0.8F_1+0.2F_2+u_1.\] 共同度按行平方和计算，特殊方差为 \(1-h_j^2\)。

\begin{center}\small\setlength{\tabcolsep}{4pt}\renewcommand{\arraystretch}{1.18}\begin{tabular}{@{}lrr@{}}
\toprule
变量 & 共同度 & 特殊方差 \\
\midrule

收缩压 & 0.68 & 0.32 \\
舒张压 & 0.5 & 0.5 \\
体重 & 0.65 & 0.35 \\
腰围 & 0.58 & 0.42 \\
\bottomrule\end{tabular}\end{center}

\textbf{（2）贡献率与命名。} 两列平方和分别为1.23、1.18，除以4得到30.75\%、29.50\%；累计贡献率为60.25\%。第一因子主要反映血压，第二因子主要反映体格或体型。旋转后的列平方和是解释方差，不应一概称为原相关矩阵的特征值。

\textbf{（3）还原相关。} 由 \(R_{\mathrm{model}}=LL^T+\Psi\)，非对角元的特殊因子协方差为0，故 \[\widehat r_{12}=0.8\times0.7+0.2\times0.1=0.58.\] 残差相关为 \(0.62-0.58=0.04\)。共同度描述一个变量的方差被解释多少；还原相关描述两个变量的联系，两者不能混用。一个残差较小也不足以单独判定整个模型拟合良好。

\textbf{（4）得分系数。} 不能直接这样做。载荷用于写``变量由因子表示''的模型；Thomson法需计算得分系数。按正文的正交模型写法，\(\widehat f=L^T(LL^T+\Psi)^{-1}z\)，通常不等于 \(L^Tz\)。




\Answer{SUP-F4}{两因子模型的研究结果评议}

\textbf{（1）拟合检验。} 对给定因子数，\(H_0\) 是总体协方差矩阵可用该因子模型表示，即 \(\Sigma=LL^T+\Psi\)。单因子模型被拒绝；两因子模型未被拒绝。在这两个候选模型中，两因子模型有更充分的拟合支持，仍应检查载荷、共同度及专业解释。

\textbf{（2）旋转作用。} 旋转前贡献率为45.00\%、20.00\%；旋转后为34.17\%、30.83\%。两者累计均为 \(3.90/6=65.00\%\)。正交旋转重新分配了单个因子的解释方差，保持共同度及总解释方差不变。第二因子贡献增加不能推出总解释率增加。

\textbf{（3）结论修订。} 可写为：``在当前样本与模型条件下，尚无充分证据否定两因子模型。''\(P=0.463\) 不是模型为真的概率，也不能证明总体恰有两个因子。命名须结合旋转载荷矩阵与各指标的含义，不能只依据因子序号或贡献率大小。
\EndExercises
